\documentclass[book.tex]{subfiles}
\begin{document}

\chapter{Applying Operator Learning to Solve Differential Equations}

\label{chap:deep_onet}
\noindent\rule{11cm}{0.4pt} \\
\textbf{Prerequisites:} \autoref{chap:dft} \\
\textbf{Difficulty Level:} ** \\
\noindent\rule{11cm}{0.4pt}
\vspace{5mm}

Deep networks satisfy a powerful universal approximation capability that allows them to learn arbitrary functions. It is less known that deep networks satisfy a universal approximation theorem with respect to learning operators and not just scalar functions. Complex nonlinear partial or ordinary differential equations provide a natural source of nonlinear operators. By using the universal operator approximation capability of deep networks, it is possible to learn the operators associated with these equations. The DeepONet \cite{lu2019deeponet} approach uses two sub-networks, one which encodes $m$ input values $x_1,\dotsc, x_m$ and another encodes $n$ output location $y$. The results of the two networks are combined to predict that value of the solution at $y$.

%"""
%While it is widely known that neural networks are universal approximators of continuous functions, a less known and perhaps more powerful result is that a neural network with a single hidden layer can approximate accurately any nonlinear continuous operator. This universal approximation theorem is suggestive of the potential application of neural networks in learning nonlinear operators from data. However, the theorem guarantees only a small approximation error for a sufficient large network, and does not consider the important optimization and generalization errors. To realize this theorem in practice, we propose deep operator networks (DeepONets) to learn operators accurately and efficiently from a relatively small dataset. A DeepONet consists of two sub-networks, one for encoding the input function at a fixed number of sensors xi,i=1,…,m (branch net), and another for encoding the locations for the output functions (trunk net). We perform systematic simulations for identifying two types of operators, i.e., dynamic systems and partial differential equations, and demonstrate that DeepONet significantly reduces the generalization error compared to the fully-connected networks. We also derive theoretically the dependence of the approximation error in terms of the number of sensors (where the input function is defined) as well as the input function type, and we verify the theorem with computational results. More importantly, we observe high-order error convergence in our computational tests, namely polynomial rates (from half order to fourth order) and even exponential convergence with respect to the training dataset size.
%"""


\section{The Universal Approximation Theorem for Operators}
Let $X$ be a Banach space. Let $K_1 \subset X$ be a compact subset. Let $K_2 \subset \mathbb{R}^d$. $V$ is a compact set in $C(K_1)$. Let $u \in V$ and $y \in K_2$. Let $G$ be some operator we seek to learn (usually $G$ is the solution operator to some potentially non-linear partial differential equation). The type of $G$ is given by

\begin{align}
    G: V \to C(K_2)
\end{align}

The universal approximation theorem for operator learning then states that for any given approximation accuracy $\epsilon > 0$. there exist constants $c_i^k, \xi_{ij}^k, \theta_i^k$, $\zeta_k$ that satisfy the approximation equation

\begin{align}
    \left | G(u)(y) - \sum_{k=1}^p \sum_{i=1}^n c_i^k \sigma \left ( \sum_{j=1}^m \xi^k_{ij} u(x_j)+ \theta^k_i \right ) \sigma(w_k \cdot y + \zeta_k) \right | < \epsilon 
\end{align}

In this equation, there are $m$ input points, $n$ hidden units in the input encoding, and $p$ hidden units in the output encoding.

\section{Basic Structure of the Deep ONet}

Let $x_1,\dotsc, x_m$ be the input points. Let $u$ be some input function. Network inputs are given by 

\begin{align}
u(x_1),\dotsc, u(x_m)
\end{align}

Here, we place a constraint that positions $x_1,\dotsc,x_m$ must be fixed in the training data. The universal approximation theorem in itself doesn't tell us how to learn an operator $G$, but we can use the statement of the bound to define "branch" and "trunk" elements


\begin{align}
\textrm{trunk}_k &= \sigma(w_k \cdot y + \zeta_k) \\
\textrm{branch}_k &= \sum_{i=1}^n c^k_i \sigma \left ( \xi^k_{ij} u(x_j) + \theta^k_i \right )
\end{align}

The core idea of the Deep ONet architecture is to approximate $G$ as an inner product of branch and trunk computed elements.

\begin{align}
G(u)(y) &= \sum_{k=1}^p \textrm{branch}_k \textrm{trunk}_k
\end{align}

The Deep ONet technique can be used for both ODEs and PDEs with different choice of operator.

%\section{ODE Solution}

%\textcolor{red}{TODO}

%\section{PDE Solution}

%\textcolor{red}{TODO}

\end{document}
